% Supplementary Information. Every number is a \nm{} key.
\section*{Supplementary Note 1: proofs}

\paragraph{Proposition~\ref{prop:invisible}.} Let $R$ be exact on the completed training answers: it accepts every
positive and rejects every negative. Let $c$ hold on every positive. Then $R\wedge c$ accepts every positive, since both
conjuncts do, and rejects every negative, since $R$ does. So $R$ and $R \wedge c$ have the same training error, and
any fit that prefers fewer conditions, as the $\ell_1$ penalty and the greedy conjunction do, returns $R$.\hfill$\square$

\paragraph{Proposition~\ref{prop:stratified}.} Let the levels be $0, 1, \ldots, L$, and let the run settle level $\lambda$
at loop $\lambda + 1$. Suppose that after loop $\lambda$ every item of level below $\lambda$ holds its value in $y$ and
every completion register of those items equals its value in $y^{+}$. An item $v$ of level $\lambda$ reads only the input
and items of level below $\lambda$, so it reads exactly what it reads at the completed answer, where its check is exact;
it therefore takes its value in $y$ at loop $\lambda+1$. The registers of $v$ are computed from items of level at most
$\lambda$, which now all hold their values in $y$, so they too equal their values in $y^{+}$. Induction on $\lambda$
gives the claim: level $L$ is settled at loop $L+1$, and at loop $L$ when level $0$ is given by the input.\hfill$\square$

\paragraph{Proposition~\ref{prop:finite}.} Let $(u,v)$ be a pair of an input $x$ of any size, and let $E_1,\ldots,E_W$ be
the witness tests of the rule and of the definition. For each $E_j$ that holds at $(u,v)$, pick one witness $k_j$, and let
$x'$ be the input restricted to $\{u,v,k_1,\ldots\}$, at most $2+W$ items, which is in the class because the class is
closed under deleting items. Own comparisons of $(u,v)$ are the same in $x$ and $x'$. A test that holds in $x$ holds in
$x'$, through its kept witness. A test that fails in $x$ fails in $x'$, since $x'$ has fewer candidate witnesses and the
legs of each are unchanged. So the rule and the definition take the same values at $(u,v)$ in $x$ as in $x'$, where they
agree by hypothesis.\hfill$\square$

\paragraph{The cold run of breadth-first search.} The CLRS-30 rule (Supplementary Note~2) names $u$ as the parent of $v$
when $\mathrm{depth}_u < \mathrm{depth}_v$, there is an edge $u \to v$, $u$ is reached, $v$ is not the source, and no
neighbour $k$ of $v$ with a lower number than $u$ has $\mathrm{depth}_k < \mathrm{depth}_v$. We show that loop $t$ sets
exactly the parents of the nodes at distance $t$ from the source. Before loop $t$, suppose the reached nodes are exactly
those at distance below $t$, with their depths right. A node $v$ at distance at least $t$ is unreached, so its depth is
$\notyet$, which is later than every number; its candidates are its reached neighbours. If $v$ is at distance $t$, these
are its neighbours at distance $t-1$ (a neighbour cannot be nearer than $t-1$), and the witness test leaves the
lowest-numbered of them, which is the library's parent. If $v$ is at distance above $t$, it has no reached neighbour, so
no parent passes; its own rule, ``no neighbour below me'', may pass instead and write the self-pointer, ``none''. That
leaves its depth at $\notyet$, since depths count from the source only, and gives way as soon as a reached neighbour
appears, because the self-pointer's rule then fails. A reached node $v$ at distance $d < t$ has candidates of depth below $d$ among
its neighbours, which lie at distance $d-1$, and the rule names the same parent as before. So after loop $t$ the reached
nodes are those at distance at most $t$, and the run comes to rest one loop after the eccentricity of the source.
The run reads every node at every loop; what keeps the far nodes from being given a parent early is the level condition:
$\mathrm{depth}_u < \mathrm{depth}_v$ is false whenever $u$ is unreached, which also makes ``$u$ reached'' redundant.

\paragraph{The greedy tasks.} For activity selection and task scheduling the answer is read only downward along one
order (finish time; weight), so the rule of item $i$ reads only items earlier in that order. The rule is then stratified
by that order, item by item, and the synchronous run fixes the first $t$ items of the order by loop $t$; it settles at
its unique fixed point, which is the greedy answer at the training answers and, by the same induction, at every input.
The string matchers are read downward along the text in the same way.

\paragraph{Proposition~\ref{prop:bf}.} Write $d^t$ for the distances computed along the pointer $\pi^t$, with $d^t_v =
\notyet$ for a node whose pointer chain does not reach the source; the source's distance is always $0$. \emph{(i) No
cycles.} Every pointer that loop $t+1$ writes satisfies $d^t_u < d^t_v$, a strict order on the numbers and $\notyet$ (it
is false between two $\notyet$), so a cycle in $\pi^{t+1}$ would need $d^t_{u_1} < \cdots < d^t_{u_1}$. \emph{(ii)
Reached nodes stay reached.} Let $v$ be reached at loop $t$, with parent $p$. Then $d^t_p = d^t_v - w(p,v) < d^t_v$,
since $w > 0$, so the set of in-neighbours $k$ with $d^t_k < d^t_v$ is not empty, and the one with the smallest
$S(k,v)$ passes: no witness has a strictly smaller $S$. The self-pointer's rule, ``no in-neighbour below me'', fails. By
induction along $d^t$, the new parent of $v$ is reached at loop $t+1$ and has a distance no larger than at loop $t$, so
$d^{t+1}_v \leq d^{t}_v$. \emph{(iii) Relaxation.} The new parent minimises $d^t_k + w(k,v)$ over the in-neighbours
below $v$, which loses no improving neighbour, since $d^t_k + w(k,v) < d^t_v$ implies $d^t_k < d^t_v$; and $d^{t+1}_v$
is the new parent's own $d^{t+1}$ plus the weight, no larger. So $d^{t+1}_v \leq \min(d^t_v, \min_k d^t_k + w(k,v))$,
and by induction $d^t_v$ is at most the length of the shortest path to $v$ with at most $t$ edges. Every $d^t_v$ is the
length of a path, so it is never below the true distance. After $h$ loops every distance is exact; loop $h+1$ writes, for
every $v$, the in-neighbour with the smallest $d_k + w(k,v)$, the shortest-path parent; and the next loop changes
nothing. The library draws positive weights: at least about $0.03$ on undirected graphs, and uniform on $(0,1)$, positive
almost surely, on directed ones. Floyd--Warshall runs this rule once per source.

\paragraph{Other cold runs.} \emph{Prim} keeps only tree edges, those whose weight equals the bottleneck, and is
breadth-first search on them from the root: loop $t$ places the nodes $t$ tree edges from the root. \emph{Depth-first
search} is fitted at every prefix of the answer's preorder and is exact there; the run places one node per loop, so
induction on the prefix gives the answer. \emph{Topological sort} reads the finish order of that run, and \emph{strong
components} run it twice, as Kosaraju's algorithm does. For \emph{longest common subsequence}, \emph{matrix chain order}
and \emph{optimal BST} the fitted rules are not stratified, and we rely on the runs: every one came to rest at the
answer on every test input.

\paragraph{The CLRS-30 certificates.} The script rebuilds each rule from its saved table and evaluates it with the code
that ran it. \emph{Sorts and minimum:} every order of distinct keys on up to \nm{cert.sorts.items} items (minimum:
\nm{cert.min.items}), against the library's labels; the library draws keys from a continuous distribution, so they are
distinct almost surely. \emph{Kruskal:} the rule reads only the pair's own conditions, the edge and the weight $w$ against
the bottleneck $B$. The axioms are that off an edge $w$ is missing, so every comparison with it is false, and that on an
edge $B \leq w$; every value they allow is checked against the cycle property, which with distinct weights says an edge
is in the tree exactly when $w = B$. Distinct weights are assumed, as continuous weights make them almost surely; the
rule also agrees with the library's Kruskal, as our generator reproduces it, on \nm{cert.kruskal.graphs} random graphs.
\emph{Breadth-first search:} the depth is given, constrained only by local axioms: at most one node at depth $0$ (the
source), and an edge joins two unreached nodes or two reached nodes whose depths differ by at most one. All
\nm{cert.bfs.cases} pair cases on up to \nm{cert.bfs.items} items with depths from 0 to 6 or $\notyet$ are checked against
the library's definition stated per pair.

\paragraph{Sudoku: the kept types at every size.} A pair of cells has a type, the pattern of which of eight relations
hold: same row, same column, same band $\lfloor r/b\rfloor$, same stack $\lfloor c/b\rfloor$, same offset in the band and
in the stack ($r \bmod b$, $c \bmod b$), and the two diagonals ($d_r = d_c$, $d_r = -d_c$, with $d_r, d_c$ the row and
column differences). Along the rows a pair is in one of four states: the same row ($d_r=0$); the same offset in another
band ($d_r$ a nonzero multiple of $b$); the same band at another offset ($0<|d_r|<b$); or neither. The columns are
alike, and the type is the two states and the two diagonal bits. A diagonal bit can hold only between states whose
differences can have equal size, and a pair with both states the same offset, or both the same band, can avoid both
diagonals only if that state allows two different sizes of difference: multiples $b, 2b, \ldots, (b-1)b$, or sizes
$1,\ldots,b-1$. Both need $b \geq 3$; every other combination is realised for every $b \geq 2$. So every $b \geq 3$
realises the same \nm{sudoku.types} types, and $b=2$ realises 23, lacking exactly these two; the script confirms the
counts for $b$ from \nm{proof.b.lo} to \nm{proof.b.hi}. The kept types are exactly those with same row, same column, or
same band and same stack, so on every board the kept pairs are the row, column and box pairs, on all three fits.

\paragraph{Sudoku: the stop is sound.} The energy of a full board is the sum of the kept costs over pairs of cells that
hold the same value. Every kept cost is at least \nm{sudoku.wmin.lo}, and none is negative. The threshold is
\nm{sudoku.tau}, below the smallest cost, so a board passes exactly when no kept pair holds the same value twice, that is,
when no row, column or box repeats a value. A unit of $n$ cells that repeats none of $n$ values holds each once. So a
board that passes is a valid Sudoku, at every size.

\paragraph{Sudoku: the cold stage never removes a solution.} The kept count conditions say that each value occurs at
least once and at most once in each unit. Both hold on every valid board of every size. Each removal of the cold stage
(and of the pair rule over overlapping units) follows from one of them, so it never removes the value that a solution
of the puzzle puts in a cell. With the complete depth-first search, every puzzle that has a solution is solved given
enough moves, and a puzzle without one is reported as such when the search is exhausted.

\paragraph{Maze-Unique.} On a tree maze of the room lattice (rooms at odd--odd cells, corridors between, walls at
even--even cells), the true answer meets every kept range: a path cell has two path side-neighbours, the start and the
goal one; an off cell has at most one, since two would close a cycle; corner counts are at most two for corridors and
zero for rooms. So no removal is wrong. From the start in which every open cell may be on the path, the cold stage removes
a cell with fewer than two possible path neighbours until none is left. By induction, every cell of the true path keeps
two possible path neighbours (the start and the goal one), so none is removed. What remains is a subtree; any branch of
it off the path ends in a leaf that is neither the start nor the goal, and that leaf is removed, so what remains at rest
is the path. No size and no edge term enter the argument, and counts run only over cells
that exist, so the edge needs no special case.

\paragraph{Maze-Hard.} The reach clock assigns each open cell the loop at which the cold reach run from the goal first
arrives, which is its side-step distance to the goal, since each loop adds exactly the cells one step further. The fitted
rule steps to an open neighbour whose clock is one less, the only ``nearer'' a grid allows, and breaks ties by the learned
order of directions, which is the order the labels use. It therefore returns the label on every maze of any size or shape.

\paragraph{CRAM's Sudoku stop.} A full $n\times n$ board that is not a valid Sudoku has at least two clashing pairs,
pairs of cells in one unit that hold the same digit. If some digit appears at least $n+1$ times, it repeats in some row
and in some column, and these are two different pairs, since two cells cannot share both a row and a column. Otherwise
every digit appears exactly $n$ times. The board is invalid, so some unit holds a digit $d$ twice. The units of that type
partition the board, and $d$ appears $n$ times in all, so another unit of that type holds no $d$; its $n$ cells hold at
most $n-1$ digits, so it repeats some digit $d'$, a second pair, disjoint from the first. So a wrong board pays at least the sum of the two smallest conflict costs, at least \nm{rel.stop.pay.lo} on every seed.
Each cell holds one digit, so the most a board can collect from negative costs is half the sum over cells of the most
one digit can collect at that cell, at most \nm{rel.stop.earn.hi}, and the input costs add less than 0.001. Every wrong
board therefore has energy at least \nm{rel.stop.bound.lo}, above the stop threshold on all three seeds.

\paragraph{CRAM's maze model inside the board.} The released Maze-Unique weights prefer a cell on exactly when at least
two of its four side neighbours are on, for all \nm{rel.cert.interior} patterns of its eight neighbours, on every seed.
On a tree maze with no open cell on the outer edge, cells only switch off in the descent from the all-on start, a cell
that has switched off never regains its second neighbour, and any order of moves ends at the path after as many moves as
there are open cells off it. No board size enters.

\section*{Supplementary Note 2: the fitted rules of all thirty tasks}

Each rule is printed from its fitted table. In a pointer rule $u$ is the candidate and $v$ the target; in a
self-pointer or item rule the item is $i$; $k$ ranges over third items, $k\notin\{u,v\}$ or $k \neq i$. ``Kept'' counts
the own conditions kept because they held on every training positive, ``added'' those added by the greedy conjunction.
A kept set may list a comparison beside its negated complement (``$\mathrm{key}_u<\mathrm{key}_v$'' and ``not
$\mathrm{key}_u\geq\mathrm{key}_v$''), because the two differ on $\notyet$, which fails every comparison. The witness
template fills both legs; a leg repeated in one test is printed once. Registers and relations: key, the input number;
target, binary search's target; pos, the item's index; x and y, a point's coordinates; s and f, an activity's start and
finish; d, a deadline (task scheduling) or a discovery time (depth-first search), and f also a finish time there; w and
edge, the weights and edges; source; depth and dist, completions, with has\_depth and has\_dist for ``reached''; $B$, the
bottleneck; $S$, the supplied sum; loop, Floyd--Warshall's self-loop weight; sel, the item mask of the partial answer;
string (an item of the text), first (the pattern's first letter) and run (the closure of agreeing text); same\_key, Uc
and Lc for longest common subsequence; $N$ and valid, the supplied split cost and ``a split inside the interval''; sep
and reach\_without, the supplied deletion closures; $M$, $m$ and Pp ($=P'$), the supplied prefix quantities, and end,
the chosen end item; open and unplaced, the supplied depth-first bits; left, right and noleft, the supplied orientation
relations, and other, the other segment's points.

{\small
% Generated by gen/make_rules.py from code/clrs/RULES.md. Do not edit by hand.
\medskip\noindent\textit{Comparisons and counts}\par
\paragraph{Insertion sort.}
\emph{pointer $u$ of $v$ ($u \neq v$)} (2 kept, 1 added; exact on the training answers):\\
\quad $\mathrm{key}_{u} < \mathrm{key}_{v}$;\\
\quad $\text{ not } \mathrm{key}_{u}  \geq  \mathrm{key}_{v}$;\\
\quad $\text{ not } \text{ some } k: (\mathrm{key}_{u} < \mathrm{key}_{k}) \text{ and } (\mathrm{key}_{k} < \mathrm{key}_{v})$;\\
\emph{$v$ points to itself (none)} (0 kept, 1 added; exact on the training answers):\\
\quad $\text{ not } \text{ some } k: (\mathrm{key}_{i}  \geq  \mathrm{key}_{k}) \text{ and } (\mathrm{key}_{k} < \mathrm{key}_{i})$;

\paragraph{Bubble sort.}
\emph{pointer $u$ of $v$ ($u \neq v$)} (2 kept, 1 added; exact on the training answers):\\
\quad $\mathrm{key}_{u} < \mathrm{key}_{v}$;\\
\quad $\text{ not } \mathrm{key}_{u}  \geq  \mathrm{key}_{v}$;\\
\quad $\text{ not } \text{ some } k: (\mathrm{key}_{u} < \mathrm{key}_{k}) \text{ and } (\mathrm{key}_{k} < \mathrm{key}_{v})$;\\
\emph{$v$ points to itself (none)} (0 kept, 1 added; exact on the training answers):\\
\quad $\text{ not } \text{ some } k: (\mathrm{key}_{i}  \geq  \mathrm{key}_{k}) \text{ and } (\mathrm{key}_{k} < \mathrm{key}_{i})$;

\paragraph{Heapsort.}
\emph{pointer $u$ of $v$ ($u \neq v$)} (2 kept, 1 added; exact on the training answers):\\
\quad $\mathrm{key}_{u} < \mathrm{key}_{v}$;\\
\quad $\text{ not } \mathrm{key}_{u}  \geq  \mathrm{key}_{v}$;\\
\quad $\text{ not } \text{ some } k: (\mathrm{key}_{u} < \mathrm{key}_{k}) \text{ and } (\mathrm{key}_{k} < \mathrm{key}_{v})$;\\
\emph{$v$ points to itself (none)} (0 kept, 1 added; exact on the training answers):\\
\quad $\text{ not } \text{ some } k: (\mathrm{key}_{i}  \geq  \mathrm{key}_{k}) \text{ and } (\mathrm{key}_{k} < \mathrm{key}_{i})$;

\paragraph{Quicksort.}
\emph{pointer $u$ of $v$ ($u \neq v$)} (2 kept, 1 added; exact on the training answers):\\
\quad $\mathrm{key}_{u} < \mathrm{key}_{v}$;\\
\quad $\text{ not } \mathrm{key}_{u}  \geq  \mathrm{key}_{v}$;\\
\quad $\text{ not } \text{ some } k: (\mathrm{key}_{u} < \mathrm{key}_{k}) \text{ and } (\mathrm{key}_{k} < \mathrm{key}_{v})$;\\
\emph{$v$ points to itself (none)} (0 kept, 1 added; exact on the training answers):\\
\quad $\text{ not } \text{ some } k: (\mathrm{key}_{i}  \geq  \mathrm{key}_{k}) \text{ and } (\mathrm{key}_{k} < \mathrm{key}_{i})$;

\paragraph{Minimum.}
\emph{the chosen item} (0 kept, 1 added; exact on the training answers):\\
\quad $\text{ not } \text{ some } k: (\mathrm{key}_{i}  \geq  \mathrm{key}_{k}) \text{ and } (\mathrm{key}_{k} < \mathrm{key}_{i})$;

\paragraph{Binary search.}
\emph{the chosen item} (0 kept, 3 added; not exact on the training answers; the score decides):\\
\quad $\#\{k: \mathrm{key}_{k} > \mathrm{key}_{i} \text{ and } \mathrm{key}_{k} < \mathrm{target}\} = \#\{k: \mathrm{key}_{k} < \mathrm{key}_{i} \text{ and } \mathrm{key}_{k}  \geq  \mathrm{target}\}$;\\
\quad $\text{ not } \#\{k: \mathrm{key}_{k} < \mathrm{key}_{i}\} < \#\{k: k = i \text{ and } \mathrm{key}_{k} < \mathrm{target}\}$;\\
\quad $\text{ not } \#\{k: k = i \text{ and } \mathrm{key}_{k} < \mathrm{target}\} = \#\{k: \mathrm{key}_{k}  \geq  \mathrm{key}_{i} \text{ and } \mathrm{key}_{k}  \geq  \mathrm{target}\}$;\\
score: +5.69\,[$\#\{k: k = i\} = \#\{k: \mathrm{key}_{k}  \leq  \mathrm{key}_{i} \text{ and } \mathrm{key}_{k}  \geq  \mathrm{target}\}$]; +3.15\,[$\#\{k: \mathrm{key}_{k} > \mathrm{key}_{i} \text{ and } \mathrm{key}_{k} < \mathrm{target}\} = \#\{k: \mathrm{key}_{k} < \mathrm{key}_{i} \text{ and } \mathrm{key}_{k}  \geq  \mathrm{target}\}$]; +1.04\,[$\#\{k: \mathrm{key}_{k}  \geq  \mathrm{key}_{i}\} = \#\{k: k = i \text{ and } \mathrm{key}_{k} < \mathrm{target}\}$]; -1.03\,[$\#\{k: k = i \text{ and } \mathrm{key}_{k} < \mathrm{target}\} < \#\{k: \mathrm{key}_{k}  \geq  \mathrm{key}_{i}\}$]; -0.90\,[$\#\{k: k = i \text{ and } \mathrm{key}_{k}  \geq  \mathrm{target}\} = \#\{k: \mathrm{key}_{k} < \mathrm{key}_{i} \text{ and } \mathrm{key}_{k}  \geq  \mathrm{target}\}$]; +0.88\,[$\#\{k: \mathrm{key}_{k} > \mathrm{key}_{i}\} = \#\{k: \mathrm{key}_{k}  \leq  \mathrm{key}_{i} \text{ and } \mathrm{key}_{k}  \geq  \mathrm{target}\}$]; +0.76\,[$\#\{k: \mathrm{key}_{k} > \mathrm{key}_{i}\} = \#\{k: \mathrm{key}_{k} < \mathrm{key}_{i} \text{ and } \mathrm{key}_{k}  \geq  \mathrm{target}\}$]; +0.64\,[$\#\{k: \mathrm{key}_{k} > \mathrm{key}_{i} \text{ and } \mathrm{key}_{k} < \mathrm{target}\} = \#\{k: \mathrm{key}_{k}  \geq  \mathrm{key}_{i} \text{ and } \mathrm{key}_{k}  \geq  \mathrm{target}\}$].

\paragraph{Quickselect.}
\emph{the chosen item} (0 kept, 2 added; exact on the training answers):\\
\quad $\text{ not } \#\{k: \mathrm{key}_{k} < \mathrm{key}_{i}\} < \#\{k: \mathrm{key}_{k} > \mathrm{key}_{i}\}$;\\
\quad $\text{ not } \#\{k: \mathrm{key}_{k}  \geq  \mathrm{key}_{i}\} < \#\{k: \mathrm{key}_{k} < \mathrm{key}_{i}\}$;

\paragraph{Activity selection.}
\emph{item selected} (0 kept, 1 added; exact on the training answers, answer read downward along f):\\
\quad $\text{ not } \text{ some } k: (\mathrm{s}_{i} < \mathrm{f}_{k}) \text{ and } (\mathrm{f}_{k} < \mathrm{f}_{i} \text{ and } \mathrm{sel}_{k})$;

\paragraph{Task scheduling.}
\emph{item selected} (0 kept, 1 added; exact on the training answers, answer read downward along w):\\
\quad $\#\{k: \mathrm{w}_{k} > \mathrm{w}_{i} \text{ and } \mathrm{sel}_{k}\} < \#\{k: \mathrm{pos}_{k} < \mathrm{d}_{i}\}$;

\paragraph{Breadth-first search.}
\emph{pointer $u$ of $v$ ($u \neq v$)} (4 kept, 1 added; exact on the training answers):\\
\quad $\mathrm{depth}_{u} < \mathrm{depth}_{v}$;\\
\quad $\mathrm{edge}(u, v)$;\\
\quad $\mathrm{has\_depth}_{u}$;\\
\quad $\text{ not } \mathrm{source}_{v}$;\\
\quad $\text{ not } \text{ some } k: (\mathrm{pos}_{u}  \geq  \mathrm{pos}_{k}) \text{ and } (\mathrm{edge}(k, v) \text{ and } \mathrm{depth}_{k} < \mathrm{depth}_{v})$;\\
\emph{$v$ points to itself (none)} (0 kept, 1 added; exact on the training answers):\\
\quad $\text{ not } \text{ some } k: (\mathrm{depth}_{i} > \mathrm{depth}_{k}) \text{ and } (\mathrm{edge}(k, i))$;

\paragraph{Kruskal.}
\emph{pair in the answer} (3 kept, 0 added; exact on the training answers):\\
\quad $\mathrm{edge}(u, v)$;\\
\quad $\mathrm{w}(u, v) = \mathrm{B}(u, v)$;\\
\quad $\text{ not } \mathrm{w}(u, v) > \mathrm{B}(u, v)$;

\paragraph{Prim.}
\emph{pointer $u$ of $v$ ($u \neq v$)} (6 kept, 0 added; exact on the training answers):\\
\quad $\mathrm{depth}_{u} < \mathrm{depth}_{v}$;\\
\quad $\mathrm{edge}(u, v)$;\\
\quad $\mathrm{w}(u, v) = \mathrm{B}(u, v)$;\\
\quad $\text{ not } \mathrm{w}(u, v) > \mathrm{B}(u, v)$;\\
\quad $\mathrm{has\_depth}_{u}$;\\
\quad $\text{ not } \mathrm{source}_{v}$;\\
\emph{$v$ points to itself (none)} (0 kept, 1 added; exact on the training answers):\\
\quad $\text{ not } \text{ some } k: (\mathrm{depth}_{i} > \mathrm{depth}_{k}) \text{ and } (\mathrm{edge}(k, i))$;

\paragraph{Longest common subsequence.}
\emph{arrow = match} (1 kept, 0 added; exact on the training answers):\\
\quad $\mathrm{same\_key}(u, v)$;\\
\emph{arrow = up} (2 kept, 0 added; exact on the training answers):\\
\quad $\text{ not } \mathrm{same\_key}(u, v)$;\\
\quad $\text{ not } \mathrm{Uc}(u, v) < \mathrm{Lc}(u, v)$;\\
\emph{arrow = left} (4 kept, 0 added; exact on the training answers):\\
\quad $\text{ not } \mathrm{same\_key}(u, v)$;\\
\quad $\mathrm{Uc}(u, v) < \mathrm{Lc}(u, v)$;\\
\quad $\text{ not } \mathrm{Uc}(u, v) = \mathrm{Lc}(u, v)$;\\
\quad $\text{ not } \mathrm{Uc}(u, v) > \mathrm{Lc}(u, v)$;

\paragraph{Naive string matcher.}
\emph{item selected} (1 kept, 1 added; exact on the training answers, answer read downward along pos):\\
\quad $\text{ not } \mathrm{string}_{i}$;\\
\quad $\#\{k: \mathrm{pos}_{k} < \mathrm{pos}_{i} \text{ and } \mathrm{sel}_{k}\} < \#\{k: \mathrm{run}(i, k) \text{ and } \mathrm{first}_{k}\}$;

\paragraph{KMP matcher.}
\emph{item selected} (1 kept, 1 added; exact on the training answers, answer read downward along pos):\\
\quad $\text{ not } \mathrm{string}_{i}$;\\
\quad $\#\{k: \mathrm{pos}_{k} < \mathrm{pos}_{i} \text{ and } \mathrm{sel}_{k}\} < \#\{k: \mathrm{run}(i, k) \text{ and } \mathrm{first}_{k}\}$;

\medskip\noindent\textit{One arithmetic primitive supplied}\par
\paragraph{Bellman--Ford.}
\emph{pointer $u$ of $v$ ($u \neq v$)} (4 kept, 1 added; exact on the training answers):\\
\quad $\mathrm{dist}_{u} < \mathrm{dist}_{v}$;\\
\quad $\mathrm{edge}(u, v)$;\\
\quad $\mathrm{has\_dist}_{u}$;\\
\quad $\text{ not } \mathrm{source}_{v}$;\\
\quad $\text{ not } \text{ some } k: (\mathrm{dist}_{k} < \mathrm{dist}_{v}) \text{ and } \mathrm{S}(k, v) < \mathrm{S}(u, v)$;\\
\emph{$v$ points to itself (none)} (0 kept, 1 added; exact on the training answers):\\
\quad $\text{ not } \text{ some } k: (\mathrm{dist}_{i} > \mathrm{dist}_{k}) \text{ and } (\mathrm{edge}(k, i))$;

\paragraph{Dijkstra.}
\emph{pointer $u$ of $v$ ($u \neq v$)} (4 kept, 1 added; exact on the training answers):\\
\quad $\mathrm{dist}_{u} < \mathrm{dist}_{v}$;\\
\quad $\mathrm{edge}(u, v)$;\\
\quad $\mathrm{has\_dist}_{u}$;\\
\quad $\text{ not } \mathrm{source}_{v}$;\\
\quad $\text{ not } \text{ some } k: (\mathrm{dist}_{k} < \mathrm{dist}_{v}) \text{ and } \mathrm{S}(k, v) < \mathrm{S}(u, v)$;\\
\emph{$v$ points to itself (none)} (0 kept, 1 added; exact on the training answers):\\
\quad $\text{ not } \text{ some } k: (\mathrm{dist}_{i} > \mathrm{dist}_{k}) \text{ and } (\mathrm{edge}(k, i))$;

\paragraph{Shortest paths in an acyclic graph.}
\emph{pointer $u$ of $v$ ($u \neq v$)} (5 kept, 1 added; exact on the training answers):\\
\quad $\mathrm{dist}_{u} < \mathrm{dist}_{v}$;\\
\quad $\mathrm{edge}(u, v)$;\\
\quad $\text{ not } \mathrm{edge}(v, u)$;\\
\quad $\mathrm{has\_dist}_{u}$;\\
\quad $\text{ not } \mathrm{source}_{v}$;\\
\quad $\text{ not } \text{ some } k: (\mathrm{dist}_{k} < \mathrm{dist}_{v}) \text{ and } \mathrm{S}(k, v) < \mathrm{S}(u, v)$;\\
\emph{$v$ points to itself (none)} (0 kept, 1 added; exact on the training answers):\\
\quad $\text{ not } \text{ some } k: (\mathrm{dist}_{i} > \mathrm{dist}_{k}) \text{ and } (\mathrm{edge}(k, i))$;

\paragraph{Floyd--Warshall.}
\emph{pointer $u$ of $v$ ($u \neq v$)} (4 kept, 1 added; exact on the training answers):\\
\quad $\mathrm{dist}_{u} < \mathrm{dist}_{v}$;\\
\quad $\mathrm{edge}(u, v)$;\\
\quad $\mathrm{has\_dist}_{u}$;\\
\quad $\text{ not } \mathrm{source}_{v}$;\\
\quad $\text{ not } \text{ some } k: (\mathrm{dist}_{k} < \mathrm{dist}_{v}) \text{ and } \mathrm{S}(k, v) < \mathrm{S}(u, v)$;\\
\emph{$v$ points to itself (none)} (0 kept, 1 added; exact on the training answers):\\
\quad $\text{ not } \text{ some } k: (\mathrm{dist}_{i} > \mathrm{dist}_{k}) \text{ and } (\mathrm{edge}(k, i))$;\\
\emph{the source's own cell} (6 kept, 1 added; exact on the training answers):\\
\quad $\mathrm{edge}(u, v)$;\\
\quad $\mathrm{w}(u, v) < \mathrm{loop}_{v}$;\\
\quad $\text{ not } \mathrm{w}(u, v) > \mathrm{loop}_{v}$;\\
\quad $\mathrm{S}(u, v) < \mathrm{loop}_{v}$;\\
\quad $\text{ not } \mathrm{S}(u, v) > \mathrm{loop}_{v}$;\\
\quad $\mathrm{has\_dist}_{u}$;\\
\quad $\text{ not } \text{ some } k: (\mathrm{dist}_{u}  \geq  \mathrm{dist}_{k}) \text{ and } (\mathrm{edge}(k, v))$;

\paragraph{Graham scan.}
\emph{item selected} (0 kept, 1 added; exact on the training answers):\\
\quad $\text{ some } k: (\mathrm{noleft}(i, k)) \text{ and } (\mathrm{noleft}(i, k))$;

\paragraph{Jarvis march.}
\emph{item selected} (0 kept, 1 added; exact on the training answers):\\
\quad $\text{ some } k: (\mathrm{noleft}(i, k)) \text{ and } (\mathrm{noleft}(i, k))$;

\paragraph{Segment intersection.}
\emph{every item passes (read-out: yes iff all pass)} (0 kept, 2 added; not exact on the training answers; the score decides):\\
\quad $\#\{k: k = i\} = \#\{k: \mathrm{left}(i, k)\}$;\\
\quad $\text{ not } \text{ some } k: (\mathrm{left}(i, k)) \text{ and } (\mathrm{left}(k, i))$;\\
score: +3.78\,[$\#\{k: k = i\} = \#\{k: \mathrm{left}(i, k)\}$]; +2.14\,[$\text{ some } k: (\mathrm{left}(i, k)) \text{ and } (\mathrm{right}(k, i))$]; +2.12\,[$\text{ some } k: (\mathrm{left}(k, i)) \text{ and } (\mathrm{right}(i, k))$]; -2.11\,[$\text{ some } k: (\mathrm{right}(i, k)) \text{ and } (\mathrm{right}(k, i))$]; -2.09\,[$\text{ some } k: (\mathrm{left}(i, k)) \text{ and } (\mathrm{left}(k, i))$]; -0.96\,[$\#\{k: \mathrm{pos}_{k}  \geq  \mathrm{pos}_{i}\} < \#\{k: \mathrm{y}_{k} > \mathrm{y}_{i}\}$]; +0.73\,[$\text{ some } k: (\mathrm{x}_{i}  \geq  \mathrm{x}_{k}) \text{ and } (\mathrm{right}(i, k))$]; +0.68\,[$\text{ some } k: (\mathrm{y}_{i} < \mathrm{y}_{k}) \text{ and } (\mathrm{other}(k, i))$].

\paragraph{Matrix chain order.}
\emph{pointer $u$ of $v$ ($u \neq v$)} (3 kept, 1 added; exact on the training answers):\\
\quad $\mathrm{pos}_{u} < \mathrm{pos}_{v}$;\\
\quad $\text{ not } \mathrm{pos}_{u}  \geq  \mathrm{pos}_{v}$;\\
\quad $\mathrm{valid}(u, v)$;\\
\quad $\text{ not } \text{ some } k: (\mathrm{pos}_{k} < \mathrm{pos}_{v}) \text{ and } \mathrm{N}(k, v) < \mathrm{N}(u, v)$;

\paragraph{Optimal search tree.}
\emph{pointer $u$ of $v$ ($u \neq v$)} (3 kept, 1 added; exact on the training answers):\\
\quad $\mathrm{pos}_{u} < \mathrm{pos}_{v}$;\\
\quad $\text{ not } \mathrm{pos}_{u}  \geq  \mathrm{pos}_{v}$;\\
\quad $\mathrm{valid}(u, v)$;\\
\quad $\text{ not } \text{ some } k: (\mathrm{pos}_{k} < \mathrm{pos}_{v}) \text{ and } \mathrm{N}(k, v) < \mathrm{N}(u, v)$;

\medskip\noindent\textit{Algorithmic structure supplied}\par
\paragraph{Maximum subarray.}
\emph{the end item} (0 kept, 1 added; exact on the training answers):\\
\quad $\text{ not } \text{ some } k: (\mathrm{M}_{i} < \mathrm{M}_{k}) \text{ and } (\mathrm{M}_{k}  \geq  \mathrm{M}_{i})$;\\
\emph{the start item (reads the end)} (0 kept, 1 added; exact on the training answers):\\
\quad $\#\{k: k = i\} = \#\{k: \mathrm{m}_{k} = \mathrm{Pp}_{i} \text{ and } \mathrm{end}_{k}\}$;

\paragraph{Articulation points.}
\emph{item selected} (0 kept, 1 added; exact on the training answers):\\
\quad $\text{ some } k: (\mathrm{edge}(i, k)) \text{ and } (\mathrm{sep}(i, k))$;

\paragraph{Bridges.}
\emph{pair in the answer} (2 kept, 0 added; exact on the training answers):\\
\quad $\mathrm{edge}(u, v)$;\\
\quad $\text{ not } \mathrm{reach\_without}(u, v)$;

\paragraph{Depth-first search.}
\emph{pointer $u$ of $v$ ($u \neq v$)} (4 kept, 3 added; exact on the training answers):\\
\quad $\mathrm{d}_{u} < \mathrm{d}_{v}$;\\
\quad $\mathrm{edge}(u, v)$;\\
\quad $\mathrm{open}_{u}$;\\
\quad $\text{ not } \mathrm{unplaced}_{u}$;\\
\quad $\text{ not } \text{ some } k: (\mathrm{edge}(u, k) \text{ and } \mathrm{d}_{u} < \mathrm{d}_{k}) \text{ and } (\mathrm{edge}(k, v) \text{ and } \mathrm{pos}_{k} < \mathrm{pos}_{v})$;\\
\quad $\text{ not } \text{ some } k: (\mathrm{d}_{u} < \mathrm{d}_{k}) \text{ and } (\mathrm{d}_{k} < \mathrm{d}_{v} \text{ and } \mathrm{open}_{k})$;\\
\quad $\text{ not } \text{ some } k: (\mathrm{edge}(u, k)) \text{ and } (\mathrm{pos}_{k} < \mathrm{pos}_{v} \text{ and } \mathrm{unplaced}_{k})$;\\
\emph{$v$ points to itself (none)} (0 kept, 2 added; exact on the training answers):\\
\quad $\#\{k: \mathrm{d}_{k} < \mathrm{d}_{i} \text{ and } \mathrm{open}_{k}\} = \#\{k: \mathrm{pos}_{k} < \mathrm{pos}_{i} \text{ and } \mathrm{unplaced}_{k}\}$;\\
\quad $\text{ not } \text{ some } k: (\mathrm{pos}_{i}  \geq  \mathrm{pos}_{k}) \text{ and } (\mathrm{pos}_{k} < \mathrm{pos}_{i} \text{ and } \mathrm{unplaced}_{k})$;

\paragraph{Topological sort.}
\emph{pointer $u$ of $v$ ($u \neq v$)} (2 kept, 1 added; exact on the training answers):\\
\quad $\mathrm{f}_{u} < \mathrm{f}_{v}$;\\
\quad $\text{ not } \mathrm{f}_{u}  \geq  \mathrm{f}_{v}$;\\
\quad $\text{ not } \text{ some } k: (\mathrm{f}_{u} < \mathrm{f}_{k}) \text{ and } (\mathrm{f}_{k} < \mathrm{f}_{v})$;\\
\emph{$v$ points to itself (none)} (0 kept, 1 added; exact on the training answers):\\
\quad $\text{ not } \text{ some } k: (\mathrm{f}_{i}  \geq  \mathrm{f}_{k}) \text{ and } (\mathrm{f}_{k} < \mathrm{f}_{i})$;\\
\emph{the last finished item} (0 kept, 1 added; exact on the training answers):\\
\quad $\text{ not } \text{ some } k: (\mathrm{f}_{i} < \mathrm{f}_{k}) \text{ and } (\mathrm{f}_{k}  \geq  \mathrm{f}_{i})$;

\paragraph{Strong components.}
No task-specific rule: the DFS check run twice (Kosaraju) and a closure read-out (dfs2.py).



\medskip\noindent\textit{Shared checks (Table~\ref{tab:shared})}\par
% Generated by gen/make_rules.py from code/clrs/final/shared.json. Do not edit by hand.
\paragraph{One check for breadth-first search, Prim, Bellman--Ford, Dijkstra and DAG shortest paths, no task named.}
\emph{pointer $u$ of $v$ ($u \neq v$)} (not exact on the training answers; a learned score decides among the candidates that pass):\\
\quad $\text{not } \mathrm{source}_{v}$;\\
\quad $\text{not } \mathrm{w}(u,v) > \mathrm{B}(u,v)$;\\
\quad $\mathrm{edge}(u,v)$;\\
\quad $\text{not } \exists k: \mathrm{dist}_{k} < \mathrm{dist}_{v} \wedge \mathrm{S}(k,v) < \mathrm{S}(u,v)$;\\
\quad $\text{not } \exists k: \mathrm{depth}_{u} \geq \mathrm{depth}_{k} \wedge \mathrm{w}(k,v) = \mathrm{B}(k,v) \wedge \mathrm{depth}_{k} < \mathrm{depth}_{v}$;\\
\emph{$v$ points to itself (none)} (exact on the training answers):\\
\quad $\#\{k: \mathrm{w}(k,i) > \mathrm{B}(k,i) \wedge \mathrm{depth}_{k} < \mathrm{depth}_{i}\} = \#\{k: \mathrm{edge}(k,i) \wedge \mathrm{source}_k\}$;\\
\quad $\text{not } \#\{k: \mathrm{w}(k,i) = \mathrm{S}(k,i) \wedge \mathrm{pos}_{k} \geq \mathrm{pos}_{i}\} < \#\{k: \mathrm{edge}(k,i) \wedge \mathrm{dist}_{k} < \mathrm{dist}_{i}\}$;\\
\quad $\#\{k: \mathrm{edge}(k,i) \wedge \mathrm{depth}_{k} < \mathrm{depth}_{i}\} = \#\{k: \mathrm{edge}(k,i) \wedge \mathrm{pos}_{k} \geq \mathrm{pos}_{i} \wedge \mathrm{source}_k\}$;

\paragraph{One check for activity selection, task scheduling, Graham scan, Jarvis march and articulation points, no task named.}
\emph{item selected} (not exact on the training answers; a learned score decides among the candidates that pass):\\
\quad $\text{not } \exists k: \mathrm{s}_{i} < \mathrm{f}_{k} \wedge \mathrm{f}_{k} < \mathrm{f}_{i} \wedge \mathrm{sel}_k$;\\
\quad $\text{not } \#\{k: \mathrm{pos}_{k} < \mathrm{d}_{i}\} < \#\{k: \mathrm{w}_{k} > \mathrm{w}_{i} \wedge \mathrm{sel}_k\}$;


}

\section*{Supplementary Table S1: the held-out battery with frozen code}

{\small Score of the library's scorer on each held-out set, in percent, with the number of inputs in brackets. Val: the
official validation split. Other columns: fresh inputs from our generator at seed $9000+n$, by size in nodes or items.
$^\dagger$Text and pattern together: 80, 320, 1,280 and 5,120 items. $^\ddagger$Always four points.
$^\S$128 items. Dashes: not run, because the run or the supplied closure is too slow at that size.\par}
\begin{center}\scriptsize
% Generated by gen/make_numbers.py from code/clrs/final/summary.json.
\begin{tabular}{lrrrrrr}\toprule
Task & val & 64 & 128 / 256 & 1,024 & sparse 1,024 & 4,096 \\\midrule
\multicolumn{7}{l}{\textit{Comparisons and counts}} \\
Insertion sort & 100.00 & 100.00 (32) & 100.00 (32) & 100.00 (32) & -- & 100.00 (4) \\
Bubble sort & 100.00 & 100.00 (32) & 100.00 (32) & 100.00 (32) & -- & 100.00 (4) \\
Heapsort & 100.00 & 100.00 (32) & 100.00 (32) & 100.00 (32) & -- & 100.00 (4) \\
Quicksort & 100.00 & 100.00 (32) & 100.00 (32) & 100.00 (32) & -- & 100.00 (4) \\
Minimum & 100.00 & 100.00 (2048) & 100.00 (256) & 100.00 (256) & -- & 100.00 (64) \\
Binary search & 100.00 & 100.00 (2048) & 100.00 (256) & 100.00 (256) & -- & 100.00 (64) \\
Quickselect & 100.00 & 100.00 (2048) & 100.00 (256) & 100.00 (256) & -- & 100.00 (64) \\
Activity selection & 100.00 & 100.00 (32) & 100.00 (32) & 100.00 (32) & -- & 100.00 (8) \\
Task scheduling & 100.00 & 100.00 (32) & 100.00 (32) & 100.00 (32) & -- & 100.00 (8) \\
Breadth-first search & 100.00 & 100.00 (32) & 100.00 (16) & 100.00 (16) & 100.00 (8) & -- \\
Kruskal & 100.00 & 100.00 (32) & 100.00 (16) & 100.00 (16) & 100.00 (8) & -- \\
Prim & 100.00 & 100.00 (32) & 100.00 (16) & 100.00 (16) & 100.00 (8) & -- \\
LCS length & 100.00 & 100.00 (32) & 100.00 (32) & 100.00 (8) & -- & -- \\
Naive string matcher$^\dagger$ & 100.00 & 100.00 (2048) & 100.00 (64) & 100.00 (64) & -- & 100.00 (4) \\
KMP matcher$^\dagger$ & 100.00 & 100.00 (2048) & 100.00 (64) & 100.00 (64) & -- & 100.00 (4) \\
\multicolumn{7}{l}{\textit{One arithmetic primitive supplied}} \\
Bellman--Ford & 100.00 & 100.00 (32) & 100.00 (16) & 100.00 (16) & 100.00 (8) & -- \\
Dijkstra & 100.00 & 100.00 (32) & 100.00 (16) & 100.00 (16) & 100.00 (8) & -- \\
DAG shortest paths & 100.00 & 100.00 (32) & 100.00 (16) & 100.00 (16) & 100.00 (8) & -- \\
Floyd--Warshall & 100.00 & 100.00 (8) & 100.00 (4) & -- & -- & -- \\
Graham scan & 100.00 & 100.00 (32) & 100.00 (32) & 100.00 (8) & -- & -- \\
Jarvis march & 100.00 & 100.00 (32) & 100.00 (32) & 100.00 (8) & -- & -- \\
Segments intersect$^\ddagger$ & 100.00 & 100.00 (2048) & -- & -- & -- & -- \\
Matrix chain order & 100.00 & 100.00 (32) & 100.00 (4)$^\S$ & -- & -- & -- \\
Optimal BST & 100.00 & 100.00 (32) & 100.00 (4)$^\S$ & -- & -- & -- \\
\multicolumn{7}{l}{\textit{Algorithmic structure supplied}} \\
Maximum subarray & 100.00 & 100.00 (1024) & 100.00 (256) & 100.00 (64) & -- & 100.00 (16) \\
Articulation points & 100.00 & 100.00 (32) & 100.00 (16) & 100.00 (4) & -- & -- \\
Bridges & 100.00 & 100.00 (32) & 100.00 (16) & -- & -- & -- \\
Depth-first search & 100.00 & 100.00 (32) & 100.00 (4) & -- & -- & -- \\
Topological sort & 100.00 & 100.00 (32) & 100.00 (4) & -- & -- & -- \\
Strong components & 100.00 & 100.00 (32) & 100.00 (4) & -- & -- & -- \\
\bottomrule\end{tabular}

\end{center}
